\section{Recomendaciones}

Las recomendaciones se ordenan por aquello que las origina: primero las que
habilitan el funcionamiento del sistema como un todo, luego las que atacan la
causa de la brecha de reconocimiento, después las que completan la medición
pendiente y, finalmente, las que desbordan el alcance de este trabajo y
constituyen líneas de continuidad.

\subsection{Validación de extremo a extremo y puesta en operación}

\begin{enumerate}
    \item \textbf{Ejecutar la simulación de extremo a extremo con las dos
        cámaras.} La cadena técnica está integrada y probada por componente, pero
        el sistema no fue ejercitado sobre una secuencia completa de ingresos y
        salidas con las dos fuentes sincronizadas. Esa corrida es la que sostiene
        la medición de $H_{1{.}5}$ y la que permite afirmar que el conjunto opera
        como un todo continuo. La infraestructura de simulación ya existe y opera
        sobre el material grabado; lo pendiente es su ejecución con registro de
        resultados.

    \item \textbf{Validar en campo contra la infraestructura del recinto.} La
        simulación no reproduce la negociación con el grabador real ni el
        comportamiento de la red local. Un despliegue piloto en un punto de
        acceso, con supervisión y sin sustituir el control manual, es la
        instancia que expone las diferencias entre el laboratorio y la operación.

    \item \textbf{Completar el banco comparativo de latencia y tasa de
        fotogramas.} El sistema ya registra la latencia de inferencia por
        fotograma, de modo que el dato se acumula desde la primera corrida. Lo
        que resta es la comparación controlada entre formatos de modelo y
        variantes de escala sobre el hardware de destino, con la degradación de
        exactitud atribuible a la media precisión.

    \item \textbf{Calibrar el umbral de confianza del emparejamiento.} El valor
        de $0{.}6$ se adoptó como política de rechazo y no como resultado de una
        curva. Calibrarlo sobre el conjunto de prueba congelado permite fijarlo
        con una tasa de error conocida, y revisar el criterio de abstención con
        datos en lugar de con criterio.

    \item \textbf{Cubrir la ruta de captura con pruebas automatizadas.} La
        verificación alcanza el dominio, los recursos y el ciclo de persistencia.
        El lector RTSP y el muestreo de fotogramas operan como proceso separado y
        no disponen de pruebas propias: es la única pieza del camino de escritura
        que se apoya en verificación manual.

    \item \textbf{Revalidar la valoración económica.} La estimación quedó cerrada
        con parámetros declarados ---$62$ puntos de función, $12$ horas por punto
        y una tarifa de Bs $30$ por hora--- que arrojan Bs 33,170. Los tres
        valores son de mercado: conviene reemplazarlos por la cotización vigente
        al momento de la presentación, sin rehacer el análisis.
\end{enumerate}

\subsection{Condiciones de captura}

Es la recomendación de mayor impacto y la de menor costo relativo, porque no
requiere modificar el software. La evidencia de la
Subsección~\ref{subsec:brecha_ocr} fija una condición concreta: para que el
reconocimiento alcance el umbral comprometido, la placa debe ocupar al menos $40$
píxeles de altura en el punto de lectura. Por debajo de ese valor el sistema
opera en un régimen donde la información se pierde antes de llegar al modelo.

En términos de instalación, esa condición se traduce en cuatro decisiones:

\begin{itemize}
    \item \textbf{Cámara dedicada al punto de acceso}, en lugar de reutilizar la
        toma general del grabador. La distancia y el ángulo deben elegirse para
        que la placa ocupe la altura mínima indicada en el instante de lectura,
        no en el mejor fotograma de la trayectoria.
    \item \textbf{Iluminación controlada en el punto de lectura.} Las capturas
        disponibles son nocturnas, con luminancia mediana de $76{.}6$ y
        desviación estándar de $20{.}7$; esa dispersión es la que obliga a
        normalizar contraste y la que degrada la lectura. Una fuente de
        iluminación fija, con filtro infrarrojo si la cámara lo admite, reduce la
        variabilidad de entrada.
    \item \textbf{Exposición y balance de blancos fijos.} El mismo vehículo gris
        claro aparece casi blanco bajo iluminación directa y casi negro a la
        sombra. Sin una referencia de exposición absoluta, la separación entre
        blanco, plata, gris y negro ---que concentra todo el error del
        clasificador de color--- no es resoluble por el modelo.
    \item \textbf{Evitar el ángulo oblicuo pronunciado.} La deformación
        perspectiva del recorte reduce la altura efectiva de los glifos y
        traslada el problema a la etapa de reconocimiento, donde ya no puede
        corregirse.
\end{itemize}

\subsection{Reconocimiento de matrículas}

Las acciones se ordenan por relación entre esfuerzo y efecto esperado, según lo
medido:

\begin{enumerate}
    \item \textbf{Escalar el conjunto de datos reales.} Es la intervención con
        mayor efecto potencial y la única que ataca la causa de fondo. La
        infraestructura de extracción y los protocolos de evaluación ya están
        construidos para escalar sin modificar el código; la tarea pendiente es
        de recolección y verificación, no de desarrollo. El objetivo debería
        fijarse en un orden de magnitud superior al volumen actual, con
        verificación humana mantenida en cada lote.

    \item \textbf{Votar sobre múltiples fotogramas de la trayectoria.} Hoy solo
        el fotograma mejor puntuado de cada vehículo alcanza la etapa de
        reconocimiento. El seguimiento ya produce decenas de observaciones por
        trayectoria; consolidarlas por votación por posición o por consenso de
        secuencia aprovecha información que el sistema ya captura y descarta.
        Esta es la técnica con mejor relación entre costo de implementación y
        ganancia esperada en coincidencia exacta.

    \item \textbf{Restringir la decodificación solo donde la gramática no es
        ambigua.} El post-procesamiento global falló porque el formato boliviano
        admite letras en el bloque numérico. Una restricción aplicada únicamente
        a las posiciones que el formato fija como numéricas, y desactivada en las
        demás, evita el efecto contrario que se midió.

    \item \textbf{Calibrar la confianza y publicar la incertidumbre.} La
        coincidencia exacta de $0{.}566$ en el sistema ensamblado indica que
        todavía una fracción apreciable de las lecturas contiene al menos un
        carácter erróneo. Reportar la confianza por carácter permite que el
        operador distinga entre una lectura utilizable y una que requiere
        confirmación.

    \item \textbf{Mantener la revisión humana como parte del ciclo.} El
        pseudo-etiquetado automático fue útil como propuesta, pero solo el
        $26{.}1\%$ de sus cajas resultó utilizable sin corrección. La revisión
        humana no es una etapa transitoria: es el mecanismo que garantiza la
        calidad del conjunto y debe presupuestarse como tal.
\end{enumerate}

\subsection{Medición pendiente por muestra insuficiente}

\begin{itemize}
    \item \textbf{Determinación de eventos ($H_{1{.}5}$).} La cadena de extremo a
        extremo ya está implementada y medida sobre una hora pareada, con una
        exactitud de $0{.}889$ cuyo intervalo contiene el umbral
        (Subsección~\ref{sec:resultados_eventos}). Lo que falta no es
        infraestructura sino muestra: con una exactitud real de $0{.}95$ se
        requieren $139$ cruces, y el recinto produce entre nueve y trece por hora
        pareada. La campaña consiste en etiquetar las horas restantes con el mismo
        procedimiento, sin volver a ajustar el sistema contra ellas.

    \item \textbf{Estabilidad de la detección medida.} La detección de vehículos
        quedó medida sobre dos cámaras y un recinto (Subsección~\ref{subsec:resultados_deteccion}).
        Repetir el conjunto sobre un tercer punto de acceso, o sobre el mismo con
        otra altura de montaje, diría si el resultado se sostiene fuera de estas
        condiciones de instalación.

    \item \textbf{Rendimiento y optimización.} Las transformaciones aplicadas
        ---exportación a formatos intermedios, media precisión y compilación de
        motores--- reducen la huella del artefacto exportado en un $75\%$, pero su
        efecto sobre latencia, tasa de fotogramas y exactitud no fue medido.
        Corresponde completar el banco de pruebas por formato y por variante de
        escala sobre el hardware de despliegue, e incluir la degradación de
        exactitud atribuible a la reducción de precisión, que es el dato que
        permite decidir si el intercambio es aceptable.
\end{itemize}

\subsection{Líneas de continuidad}

\begin{enumerate}
    \item \textbf{Integración del padrón con barreras físicas y mecanismos RFID.}
        Una vez implementado el padrón lógico de vehículos autorizados (RF-10) en
        software, la línea de evolución natural consiste en la integración con
        actuadores físicos (barreras vehiculares, semáforos de acceso o control
        de apertura automática) o sistemas complementarios de radiofrecuencia (RFID).

    \item \textbf{Clasificación de color con referencia de exposición.} El error
        se concentra en el bloque acromático y no es de tono sino de nivel de
        brillo. Un modelo que reciba una referencia de exposición ---una región
        de calibración en la escena o una tarjeta de grises--- podría separar
        blanco, plata, gris y negro con un criterio distinto al actual.

    \item \textbf{Reidentificación y seguimiento entre cámaras.} La arquitectura
        prevé dos puntos de acceso. La asociación de una misma visita entre
        cámaras permitiría determinar eventos sin depender exclusivamente de la
        lectura de matrícula, lo que reduce la exposición a la brecha de
        reconocimiento.

    \item \textbf{Despliegue en hardware de borde.} Los motores compilados son
        específicos del dispositivo y del controlador con que se generaron. Un
        procedimiento reproducible de compilación por equipo de destino es
        condición para distribuir el sistema más allá del equipo de desarrollo.

    \item \textbf{Tratamiento de datos personales.} Las matrículas y los registros
        asociados constituyen datos personales. El sistema aplica una retención de
        $90$ días para la telemetría y los recortes, y conserva el vehículo y la
        sesión para la trazabilidad. Antes de una operación real debe definirse la
        base de licitud, el plazo definitivo de conservación, quién accede a la
        información y cómo se atiende una solicitud de supresión; la política de
        retención es, en último término, una decisión institucional que el software
        no resuelve por sí solo.

    \item \textbf{Extensión del conjunto de datos a otros recintos.} Todas las
        sesiones provienen de las mismas cámaras. La partición por bloque horario
        controla la fuga por identidad, pero no la variabilidad entre modelos de
        cámara o condiciones de instalación distintas. Un conjunto con material de
        un segundo recinto permitiría medir la generalización, que es la pregunta
        que un despliegue real plantea de inmediato.
\end{enumerate}

\subsection{Consideración final}

El trabajo deja un sistema construido, un conjunto de datos verificado y un
conjunto de mediciones que delimitan con precisión qué sabe hacer el sistema y
qué no. De los siete indicadores de contraste, dos se respaldan, tres quedan por
debajo del umbral o de la referencia y dos no se contrastan; de los seis objetivos
específicos, tres se cumplen y tres se cumplen parcialmente. Esa proporción no es un resultado adverso: es el
mapa exacto del estado del arte alcanzado, y es la información que permite que
cualquier trabajo posterior se apoye en lo construido sin repetir lo ya medido.
